\documentclass{article}
\usepackage[UTF8]{ctex}
\usepackage{amsmath,amssymb,amsfonts}
\usepackage{enumitem}
\usepackage{geometry}
\geometry{a4paper, margin=1in}
% === 错题本排版优化（不参与题目标识） ===
\setlength{\parindent}{0pt}
\setlength{\parskip}{0.35em plus 0.10em minus 0.08em}
\setlength{\emergencystretch}{3em}
\setlist{itemsep=0.25em, topsep=0.40em, parsep=0pt}
\allowdisplaybreaks[2]
% === 错题本排版优化结束 ===
\begin{document}
\section*{南开大学2024年数学分析真题}
\section*{55．南开大学2024年数学分析试卷}

\subsection{题目1}
1．（20 分）求极限

\[
\lim _{x \rightarrow 0}\left[\frac{e^{x}}{\ln (1+x)}-\frac{1}{x}\right]
\]

\subsection{题目2}
2．（20 分）设 \(z\) 是由方程

\[
2 x^{3} z-3 x^{2} y^{2}-2 x y z^{2}-4 y^{3} z+4=0
\]

确定的关于 \(x, y\) 的隐函数，在 \(x=2, y=1, z=2\) 处计算全微分 \(\mathrm{d} z\) 。

\subsection{题目3}
3．（20 分）判断广义积分 \(\int_{0}^{+\infty} \frac{x \ln x}{x^{3}+1} \mathrm{~d} x\) 的敛散性。

\subsection{题目4}
4．（25 分）计算曲线积分

\[
\int_{L}\left(x^{3} y+e^{z}\right) \mathrm{d} x+\left(x y^{3}+5 x^{3} y^{2}+x e^{y}\right) \mathrm{d} y
\]

其中 \(L\) 为椭圆 \(4 x^{2}+9 y^{2}=36\) ，取逆时针方向．

\subsection{题目5}
5．（25 分）求幂级数 \(\sum_{n=1}^{\infty} n(n+1)\left(\frac{x-1}{x}\right)^{n}\) 的收敛域及和函数．

\subsection{题目6}
6．（15 分）设 \(f(x)\) 在 \([0,1]\) 上单调递增，\(f(0)=0, f(1)=1, n\) 为正整数，\(\xi_{k} \in\left[\frac{k-1}{n}, \frac{k}{n}\right], k=1,2, \cdots, n\) ，证明：

\[
\left|\int_{0}^{1} f(x) \mathrm{d} x-\frac{1}{n} \sum_{k=1}^{n} f\left(\xi_{k}\right)\right| \leq \frac{1}{n}
\]

\subsection{题目7}
7．（15 分）设 \(f(x)\) 在 \([a, b]\) 上连续，\(f(a)<f(b)\) ，证明：存在 \([c, d] \subseteq[a, b]\) ，使得 \(f(x)\) 在 \([c, d]\) 上的最小值为 \(f(a)\) ，最大值为 \(f(b)\) ．

\subsection{题目8}
8．（10 分）设 \(\sum_{n=1}^{\infty} a_{n}\) 是一个收敛的正项级数，令 \(f(x)=\sum_{n=1}^{\infty} a_{n}|\sin n x|\) ，已知 \(f(x)\) 在 \((-\infty,+\infty)\) 上满足利普希茨条件，即存在常数 \(L>0\) ，使得对任意的 \(x, y \in(-\infty,+\infty)\) ，均有

\[
|f(x)-f(y)| \leq L|x-y|
\]

证明：级数 \(\sum_{n=1}^{\infty} n a_{n}\) 收敛。
\end{document}
